% ============================================================================
% COURS 08 — ÉLECTROMÉCANIQUE
% ============================================================================

\section{Conversion électromécanique de l'énergie}\label{sec:em-conversion}

L'électromécanique étudie les dispositifs qui convertissent l'énergie électrique
en énergie mécanique, ou inversement. Les moteurs, génératrices, actionneurs
linéaires et transformateurs reposent tous sur les interactions entre courants,
champs magnétiques et mouvements.

\TLChaineConversionElectromecanique

Pour un convertisseur tournant :
\[
P_e=u i,\qquad P_m=C\Omega,\qquad \eta=\frac{P_{\mathrm{utile}}}{P_{\mathrm{absorbée}}}.
\]

\section{Grandeurs magnétiques}\label{sec:em-magnetisme}

Le champ magnétique est caractérisé par le champ $\vec H$ et l'induction
magnétique $\vec B$. Dans un matériau linéaire :
\[
\vec B=\mu\vec H=\mu_0\mu_r\vec H.
\]
Le flux traversant une surface $S$ vaut
\[
\Phi=\iint_S \vec B\cdot\mathrm d\vec S.
\]

\subsection{Circuit magnétique}

\TLCircuitMagnetiqueSimple

L'analogie avec un circuit électrique conduit aux relations
\[
\mathcal F=Ni,\qquad \mathcal R=\frac{\ell}{\mu S},\qquad
\Phi=\frac{\mathcal F}{\mathcal R}.
\]
La force magnétomotrice $\mathcal F$ joue le rôle d'une tension, le flux $\Phi$
celui d'un courant et la réluctance $\mathcal R$ celui d'une résistance.

La présence d'un entrefer augmente fortement la réluctance totale :
\[
\mathcal R_{\mathrm{tot}}=\mathcal R_{\mathrm{fer}}+\mathcal R_{\mathrm{entrefer}}.
\]

\section{Induction électromagnétique}\label{sec:em-induction}

La loi de Faraday-Lenz relie la variation du flux à la force électromotrice :
\[
e(t)=-N\frac{\mathrm d\Phi}{\mathrm dt}.
\]
Le signe traduit l'opposition du phénomène induit à la cause qui lui donne
naissance.

\subsection{Bobine réelle}

Une bobine réelle possède une inductance $L$ et une résistance $R$.

\TLBobineRL

Son équation électrique est
\[
u(t)=Ri(t)+L\frac{\mathrm di(t)}{\mathrm dt}.
\]
Pour une alimentation continue par un échelon $U$, le courant s'établit selon
\[
i(t)=\frac{U}{R}\left(1-e^{-t/\tau}\right),\qquad \tau=\frac{L}{R}.
\]
L'énergie stockée dans le champ magnétique vaut
\[
E_m=\frac12Li^2.
\]

\section{Transformateur}\label{sec:em-transformateur}

Un transformateur transfère une puissance alternative entre deux circuits grâce
à un flux magnétique commun.

\TLTransformateurIdeal

Pour un transformateur idéal :
\[
\frac{u_2}{u_1}=\frac{N_2}{N_1}=m,
\qquad
\frac{i_2}{i_1}=\frac{N_1}{N_2}=\frac1m,
\qquad
P_1=P_2.
\]
En pratique, les pertes proviennent principalement des résistances des
bobinages, des courants de Foucault, de l'hystérésis et des fuites de flux.

\section{Forces et couples électromagnétiques}\label{sec:em-efforts}

Un conducteur parcouru par un courant dans un champ magnétique subit la force de
Laplace :
\[
\mathrm d\vec F=i\,\mathrm d\vec \ell\wedge\vec B.
\]
Dans une machine tournante, la somme de ces efforts produit un couple. Une forme
générale du bilan énergétique conduit à
\[
C_{em}=\frac{\partial W'_m}{\partial\theta},
\]
où $W'_m$ est la coénergie magnétique.

\section{Modèle élémentaire d'un moteur à courant continu}\label{sec:em-mcc}

Même si l'étude détaillée de la machine à courant continu fait l'objet du cours
suivant, son modèle élémentaire illustre directement la conversion
électromécanique.

\TLModeleMoteurCC

Les équations sont
\[
u=Ri+L\frac{\mathrm di}{\mathrm dt}+e,
\qquad e=K_e\Omega,
\qquad C_{em}=K_t i.
\]
L'équation mécanique s'écrit
\[
J\frac{\mathrm d\Omega}{\mathrm dt}=C_{em}-C_r-f\Omega.
\]
Dans le système international et pour une machine idéale, les constantes $K_e$
et $K_t$ ont la même valeur numérique.

\section{Bilans de puissance et rendement}\label{sec:em-puissances}

\TLBilanPuissancesElectromecanique

Pour un moteur :
\[
P_e=ui,\qquad P_{em}=C_{em}\Omega=e i,\qquad P_u=C_u\Omega.
\]
Les pertes Joule valent
\[
P_J=Ri^2.
\]
Le rendement global est
\[
\eta=\frac{P_u}{P_e}.
\]

\section{Régimes transitoires et permanents}\label{sec:em-regimes}

À court terme, l'inductance limite les variations du courant. À plus long terme,
l'inertie mécanique limite les variations de vitesse. Une chaîne
électromécanique possède donc généralement au moins deux constantes de temps :
une constante électrique et une constante mécanique.

En régime permanent :
\[
\frac{\mathrm di}{\mathrm dt}=0,
\qquad
\frac{\mathrm d\Omega}{\mathrm dt}=0.
\]
Les équations différentielles deviennent alors algébriques, ce qui facilite le
tracé des caractéristiques couple--vitesse.

\section{À retenir}\label{sec:em-retenir}

\begin{itemize}
  \item La force magnétomotrice $Ni$ crée un flux dans un circuit magnétique.
  \item La réluctance augmente fortement lorsqu'un entrefer est introduit.
  \item Une variation de flux induit une force électromotrice.
  \item Une bobine stocke l'énergie $\frac12Li^2$.
  \item Un transformateur adapte tension et courant sans changer la fréquence.
  \item Dans une machine tournante, la puissance électromagnétique vérifie
  $P_{em}=C_{em}\Omega$.
  \item Les pertes électriques, magnétiques et mécaniques limitent le rendement.
\end{itemize}
